\subsection{W(A) 的 Witt 向量环}

设 A 为一个环。若 \(\boldsymbol{a} = (a_n)_{n \in \mathbb{N}}\) 和 \(\boldsymbol{b} = (b_n)_{n \in \mathbb{N}}\) 是 \(\mathbf{A}^{\mathbf{N}}\) 的元素，则以 \(S_{\mathbf{A}}(\boldsymbol{a}, \boldsymbol{b})\)（分别以 \(P_{\mathbf{A}}(\boldsymbol{a}, \boldsymbol{b})\)，以及 \(I_{\mathbf{A}}(\boldsymbol{a})\)），或简称 \(S(\boldsymbol{a}, \boldsymbol{b})\)（分别简称 \(P(\boldsymbol{a}, \boldsymbol{b})\)，以及 \(I(\boldsymbol{a})\)）表示序列 \((S_n(a_0, ..., a_n; b_0, \ldots, b_n))_{n \in \mathbb{N}}\)（分别为 \((P_n(a_0, ..., a_n; b_0, \ldots, b_n))_{n \in \mathbb{N}}\)，以及 \((I_n(a_0, ..., a_n))_{n \in \mathbb{N}}\)）。在公式 (12)、(13) 和 (14) 中，以 \(a_n\) 代替 \(X_n\)、以 \(b_n\) 代替 \(Y_n\)，对所有 \(n \in \mathbb{N}\)，得到等式

\begin{enumerate}
\def\labelenumi{(\arabic{enumi})}
\setcounter{enumi}{15}
\tightlist
\item
\end{enumerate}

\[
\Phi_ {\mathrm{A}} (\mathrm{S} _ {\mathrm{A}} (a, b)) = \Phi_ {\mathrm{A}} (a) + \Phi_ {\mathrm{A}} (b)\tag{17}
\]

\[
\Phi_ {\mathrm{A}} (\mathrm{P} _ {\mathrm{A}} (a, b)) = \Phi_ {\mathrm{A}} (a) \Phi_ {\mathrm{A}} (b)\tag{18}
\]

\[
\Phi_ {\mathbf {A}} \left(\mathrm{I} _ {\mathbf {A}} (\boldsymbol {a})\right) = - \Phi_ {\mathbf {A}} (\boldsymbol {a}).
\]

以 W(A) 表示集合 A \(^{N}\)，并赋予复合律 S \(_{A}\) 和 P \(_{A}\)。

  设 \(\rho : \mathbf{B} \to \mathbf{A}\) 为环同态。以 \(\rho^{\mathbf{N}}\)，或仍以 \(\mathbf{W}(\rho)\)，表示从 \(\mathbf{B}^{\mathbf{N}}\) 到 \(\mathbf{A}^{\mathbf{N}}\) 的映射，它将元素 \(b = (b_n)_{n \in \mathbb{N}}\)（属于 \(\mathbf{B}^{\mathbf{N}}\)）送到 \((\rho(b_n))_{n \in \mathbb{N}}\)。由定义立即得到

\begin{enumerate}
\def\labelenumi{(\arabic{enumi})}
\setcounter{enumi}{18}
\tightlist
\item
\end{enumerate}

\[
\mathrm{W} (\rho) \circ \mathrm{S} _ {\mathrm{B}} = \mathrm{S} _ {\mathrm{A}} \circ (\mathrm{W} (\rho) \times \mathrm{W} (\rho))\tag{20}
\]

\[
\mathrm{W} (\rho) \circ \mathrm{P} _ {\mathrm{B}} = \mathrm{P} _ {\mathrm{A}} \circ (\mathrm{W} (\rho) \times \mathrm{W} (\rho))\tag{21}
\]

\[
\mathrm{W} (\rho) \circ \mathrm{I} _ {\mathrm{B}} = \mathrm{I} _ {\mathrm{A}} \circ \mathrm{W} (\rho)\tag{22}
\]

\[
\rho^ {\mathbf {N}} \circ \Phi_ {\mathrm{B}} = \Phi_ {\mathrm{A}} \circ \mathrm{W} (\rho).
\]

 引理 3. --- 设 A 为一个环。存在满射环同态 \(\rho: \mathbf{B} \to \mathbf{A}\)，其中 \(\mathbf{B}\) 为满足下列条件的环：p 不是 \(\mathbf{B}\) 中的零因子，且存在自同态 \(\sigma\)，它是 \(\mathbf{B}\) 的自同态，使得 \(\sigma(b) \equiv b^p\) mod. \(p\) . \(\mathbf{B}\) 对每个 \(b \in \mathbf{B}\) 成立。

  事实上，只需置 \(\mathbf{B} = \mathbf{Z}[(X_{a})_{a \in A}]\)，取 \(\sigma\) 为 \(\mathbf{B}\) 的自同态，由 \(\sigma(X_{a}) = X_{a}^{p}\) 对每个 \(a \in A\) 定义，并取 \(\rho\) 为从 \(\mathbf{B}\) 到 A 的同态，由 \(\rho(X_{a}) = a\) 对每个 \(a \in A\) 定义即可。

定理 1. --- a) 设 A 为一个（交换）环。赋予 \(S_A\) 加法和 \(P_A\) 乘法后，\(W(A)\) 是一个（交换）环。加法单位元是序列 \(0_A\)，其所有项均为零；乘法单位元是序列 \(1_A\)，除指标为 0 且等于 \(1_A\) 的项外，其余各项均为零。元素 \(a\)（属于 \(W(A)\)）的加法逆元是 \(I_A(a)\)。

\begin{enumerate}
\def\labelenumi{\alph{enumi})}
\setcounter{enumi}{1}
\item
  设 \(\rho: \mathbf{B} \to \mathbf{A}\) 为环同态。则 \(\mathbf{W}(\rho): \mathbf{W}(\mathbf{B}) \to \mathbf{W}(\mathbf{A})\) 是环同态。
\item
  设 A 为一个环。映射 \(\Phi_{\mathbf{A}}\) 是从 \(W(A)\) 到乘积环 \(\mathbf{A}^{\mathbf{N}}\) 的环同态。特别地，对于每个 \(n \in \mathbf{N}\)，映射 \(\Phi_{n}: a \mapsto \Phi_{n}(a_{0}, \ldots, a_{n})\) 是从 W(A) 到 A 的环同态。
\end{enumerate}

  由公式 (16)、(17)、(19) 和 (20)，只需证明断言 \(a\)。设 \(\rho: \mathbf{B} \to \mathbf{A}\) 为满足引理 3 条件的环同态。令 \(\mathbf{B}'\) 为 \(\mathbf{B}^{\mathbf{N}}\) 的子环，由元素 \((b_n)_{n \in \mathbf{N}}\) 构成，且满足 \(\sigma(b_n) \equiv b_{n+1} \pmod{p^{n+1}\mathbf{B}}\) 对所有 \(n \in \mathbf{N}\) 成立。根据第 2 号的命题 2，\(\Phi_{\mathbf{B}}\) 诱导出 \(\Phi_{\mathbf{B}}'\)，它是从 W(B) 到 \(\mathbf{B}'\) 的双射。结合公式 (16) 至 (18) 以及关系 \(\Phi_n(\mathbf{0}_{\mathbf{B}}) = 0\) 和 \(\Phi_n(\mathbf{1}_{\mathbf{B}}) = \mathbf{1}_{\mathbf{B}}\)（\(n \in \mathbf{N}\)），通过结构迁移可见 W(B) 是一个环，其加法单位元为 \(\mathbf{0}_{\mathbf{B}}\)，乘法单位元为 \(\mathbf{1}_{\mathbf{B}}\)，且 b 的相反元等于 \(\mathbf{I}_{\mathbf{B}}(b)\)。

映射 W(ρ): W(B) → W(A) 是满射。根据公式 (19) 和 (20)，与映射 W(ρ) 相关的 W(B) 上的等价关系 R 与 W(B) 的环结构相容。由于 W(ρ) 诱导出从商环 W(B)/R 到 W(A) 的双射 Ψ，且与加法和乘法相容，断言 a) 由结构迁移得出。

定义 1. --- 设 A 为一个环。称环 \(W(A)\) 为系数在 A 中的 Witt 向量环。

对于元素 \(\boldsymbol{a}\)（属于 \(W(A)\)）及 \(n\)（属于 \(\mathbf{N}\)），元素 \(\Phi_n(\boldsymbol{a}) = \Phi_n(a_0, \ldots, a_n)\) 有时称为指标 \(n\) 的 \(\boldsymbol{a}\) 的幽灵分量。

注记。--- 重新采用第 3 号注记的记号。设 A 为一个环。若 \(a\) 和 \(b\) 是 \(\mathbf{A}^{\mathbf{J}}\) 中的元素，且 \(r = (r_{j})_{j\in\mathbf{J}}\) 是 \(\mathcal{A}^{\mathbf{J}}\) 中的元素，则以 \(r_{\mathrm{A}}(a,b)\) 表示元素 \((r_{j}(a,b))_{j\in\mathbf{J}}\)，它属于 \(\mathbf{A}^{\mathbf{J}}\)。以 \(U(A)\) 表示集合 \(\mathbf{A}^{\mathbf{J}}\)，并赋予其复合律 \(s_{\mathrm{A}}\) 和 \(p_{\mathrm{A}}\)。可以证明（p. 52, 习题 35），赋予加法 \(s_{\mathrm{A}}\) 和乘法 \(p_{\mathrm{A}}\) 后，\(U(A)\) 是一个（交换）环；称其为 A 的万有 Witt 环。加法单位元是 \(U(A)\) 中所有分量均为零的元素；乘法单位元是 \(U(A)\) 中除指标为 1 且等于 \(1_{\mathrm{A}}\) 的分量外其余分量均为零的元素；元素 \(a\)（属于 \(U(A)\)）的相反元是 \(i_{\mathrm{A}}(a)\)。映射 \(\varphi_{\mathrm{A}}\) 是从 \(U(A)\) 到乘积环 \(\mathbf{A}^{\mathbf{J}}\) 的环同态。

设 \(\rho: \mathbf{B} \to \mathbf{A}\) 为环同态；以 \(\mathrm{U}(\rho)\) 表示从 \(\mathbf{B}^{\mathbf{J}}\) 到 \(\mathbf{A}^{\mathbf{J}}\) 的映射，它将元素 \((b_j)_{j \in \mathbf{J}}\)（属于 \(\mathbf{B}^{\mathbf{J}}\)）送到元素 \((\rho(b_j))_{j \in \mathbf{J}}\)（属于 \(\mathbf{A}^{\mathbf{J}}\)）。可以证明（同上）\(\mathrm{U}(\rho)\) 是从 \(\mathrm{U}(\mathrm{B})\) 到 \(\mathrm{U}(\mathrm{A})\) 的环同态。
% semantic-fragments: legacy-input-seam
\relax{}
